\section*{Вступ}
\phantomsection
\addcontentsline{toc}{section}{Вступ}

\emph{Анонімні повноваження (credential)} дозволяють користувачеві отримати від емітента сертифікат певних атрибутів 
і згодом довести, що він володіє таким сертифікатом, не розкриваючи при цьому нічого про себе, крім того, що потрібно 
верифікатору; різні представлення одних і тих самих повноважень не можуть бути пов'язані між собою. Ідея сягає 
часів Чаума (Chaum)~\cite{Chaum85}, і сьогодні вона реалізована у двох сімействах (\ref{fig:taxonomy}). У сімействі 
з \emph{публічною верифікацією} -- посвідченням є цифровий підпис, наприклад, Camenisch–Lysyanskaya (CL)~\cite{CL01} 
або BBS/BBS+~\cite{BBS04,ASM06,TZ23}, причому останній наразі стандартизовано IRTF CFRG (Internet Research Task Force 
Crypto Forum Research Group)~\cite{BBSdraft}. У сімействі з \emph{ключовою верифікацією} (keyed-verification) (KVAC) 
емітент є також єдиним верифікатором, тому підпис можна замінити алгебраїчним MAC~\cite{CMZ14}, який працює у звичайній 
групі порядку простого числа без білінійних спарювань.

На думку спада повсякденна аналогія: паспорт перевіряє будь-хто, хто знає державний печатний штамп (піддається 
перевірці публічно), тоді як жетон гардеробу перевіряє лише той гардероб, який його видав (ключова верифікація) 
-- і саме тому, що ніхто інший його не перевіряє, його можна зробити дешевшим.

\begin{figure}[!ht]
  \centering
  \begin{minipage}{0.95\linewidth}
    \centering
    % Two families of anonymous credentials and where Signal's zkgroup and BBS sit.
% Requires: \usepackage[svgnames]{xcolor}, \usetikzlibrary{positioning}
\begin{tikzpicture}[
    font=\small,
    node distance=4mm and 4mm,
    root/.style={draw=DarkOrchid, fill=Thistle, rounded corners=3pt, font=\small\bfseries, align=center,
                 minimum width=4cm},
    fam/.style={draw=DarkOrchid, fill=Thistle!45, rounded corners=3pt, align=center, text width=5.6cm},
    leaf/.style={draw=Plum, fill=white, rounded corners=2pt, align=left, font=\footnotesize, text width=5.6cm},
    hl/.style={leaf, draw=DarkOrchid, very thick, fill=Lavender},
    edge/.style={DarkOrchid, thick}
  ]
  \node[root] (R) {Анонімні повідомлення\\[-1pt]\footnotesize Chaum 1985};
  \node[fam, below left=7mm and -1.2cm of R]  (PV) {\textbf{Publicly verifiable}\\[-1pt]
        \footnotesize посвідчення $=$ цифровий підпис};
  \node[fam, below right=7mm and -1.2cm of R] (KV) {\textbf{Keyed-verification (KVAC)}\\[-1pt]
        \footnotesize посвідчення $=$ алгебраїчний MAC};
  \node[leaf, below=of PV] (pv1) {CL підписи (strong RSA, 2001)};
  \node[hl,   below=1.5mm of pv1] (pv2) {BBS / BBS+ (спарювання, $q$-SDH);\\ IETF CFRG draft, BLS12-381};
  \node[leaf, below=1.5mm of pv2] (pv3) {PS підписи, U-Prove, \dots};
  \node[leaf, below=of KV] (kv1) {CMZ: $\mathsf{MAC}_{\mathsf{GGM}}$, $\mathsf{MAC}_{\mathsf{DDH}}$ (2014)};
  \node[hl,   below=1.5mm of kv1] (kv2) {Chase-Perrin-Zaverucha (2020)\\ $\rightarrow$ Signal \texttt{zkgroup}};
  \node[leaf, below=1.5mm of kv2] (kv3) {BBS-MAC (BBDT 2016), IETF ARC, $\mu$CMZ/$\mu$BBS (2024), Tor Lox};
  \draw[edge] (R) -- (PV); \draw[edge] (R) -- (KV);
  \draw[edge] (PV) -- (pv1); \draw[edge] (KV) -- (kv1);
  \node[font=\footnotesize, align=center, below=5mm of $(pv3.south)!0.5!(kv3.south)$]
       {Загальний механізм: $\Sigma$-протоколи $+$ Фіат-Шамір для доказів знання/володіння відповідними даними};
\end{tikzpicture}

    \caption{Два сімейства анонімних посвідчень. Жирним контуром виділив систему, що досліджується (праворуч), 
      та альтернативу, проаналізовану в розділі~\ref{sec:alternatives} (ліворуч).}
    \label{fig:taxonomy}
  \end{minipage}
\end{figure}

\textbf{Чому я образ саме цю тему дослідження?} \noindent На це в мене було дві причини, окрім банальної цікавості)).
\begin{enumerate}
  \item \textbf{Актуальність.} Приватні групи Signal використовують KVAC з моменту їхнього перепроектування, 
    оголошеного у 2019 році~\cite{SignalBlog2019,CPZ20}, у додатку з мільйонами користувачів; той сами підхід 
    використовує система розподілу мостів Lox у Tor для захисту соціальної мережі своїх користувачів~\cite{Lox23}. 
    Ця галузь швидко розвивається: проект CFRG щодо $\Sigma$-протоколів~\cite{SigmaDraft}, дванадцята редакція 
    проекту BBS (вересень 2026 року)~\cite{BBSdraft}, токени з обмеженою швидкістю (rate-limited) на основі 
    KVAC~\cite{ARC25} та пропозиція 2026 року від дослідників Signal та IBM щодо квантово-безпечної реорганізації 
    системи груп~\cite{Connell26}.
  \item \textbf{Відкрита реалізація для експериментів.} libsignal є відкритим програмним забезпеченням (написаний 
    на Rust) і компілюється на звичайному ноутбуці, тому наведені цифри можна перевірити: я самостійно затестив 
    і відтворив тести продуктивності та порівняв їх із реалізацією BBS з відкритим кодом~\cite{zkryptium} (розділ~\ref{sec:perf}).
\end{enumerate}

\noindent Протягом звіту намагався дотримуватися структурі дослідження: у розділі~\ref{sec:project} описав проблему 
та роль нульового знання, у розділі~\ref{sec:protocol} -- систему доказів та її адаптацію, а в розділі~\ref{sec:analysis} 
-- припущення, відомі ризики, продуктивність та можливі альтернативи.

\section{Опис проєкту: приватні групи в Signal}
\label{sec:project}

\subsection{Проблема та попереднє рішення}
Груповий чат потребує \emph{стану}: списку учасників, ролей (адміністратор або учасник), імені, аватара, таймера 
для зникаючих повідомлень. Звичайний сервер зберігає цей стан у вигляді відкритого тексту, може забезпечувати 
контроль доступу (\enquote{лише адміністратор може додавати учасників}) -- і тим самим отримує інформацію про 
соціальний граф усіх своїх користувачів. До оновлення в 2019 Signal уникав цього, зберігаючи групи виключно на 
клієнтах: сервер не мав записів про членство, а оновлення груп надсилалися як звичайні повідомлення з end-to-end 
шифруванням~\cite{SignalBlog2019}. Ціною цього була відсутність канонічного стану групи (клієнти могли розходитися) 
та неможливість контролю доступу на основі ролей: кожен учасник мав однакові права, оскільки сервер не міг 
перевірити привілейованість~\cite{SignalBlog2019}.

\subsection{Нова архітектура та необхідність використання zero-knowledge схеми}
Система приватних груп (Private Group System, PGS) повертає стан на сервер, але у зашифрованому вигляді (рис.~\ref{fig:architecture}). 
Кожна група має випадковий \emph{GroupMasterKey}, який передається між учасниками end-to-end; на його основі 
клієнти виводять \emph{GroupSecretParams}, що використовуються для шифрування ідентифікатора ($\UID$) та ключа 
профілю ($\PK$) для кожного учасника. Сервер зберігає лише зашифровані тексти та перевіряє анонімні представлення. 
Коли Аліса хоче прочитати або змінити дані групи, вона повинна переконати сервер, що вона є \emph{одним із зашифрованих учасників} 
-- не розкриваючи, яким саме, і не даючи серверу можливості пов'язати її запити з її обліковим записом~\cite{CPZ20}. 
Це найбільш типовий приклад прояву zero-knowledge.

Як повсякденна аналогія: Аліса показує охоронцю запечатаний конверт ($E$) разом із сертифікатом ($\pi$), у якому 
зазначено, що у конверті міститься її справжнє посвідчення особи. Охоронець знаходить ідентичний конверт у списку 
гостей і впускає її навіть не відкриваючи його і не дізнаючись, чиє посвідчення знаходиться всередині. Коли запис 
про учасника додається або оновлюється, сертифікат ProfileKeyCredential додатково гарантує охоронцю, що 
другий конверт ($E'$) містить ключ профілю, який до цього був сертифікований сервером для цієї особи.

\begin{figure}[!ht]
  \centering
  \begin{minipage}{0.95\linewidth}
    \centering
    % Architecture of the Signal Private Group System (after Chase, Perrin, Zaverucha, CCS 2020).
% Requires: \usepackage[svgnames]{xcolor}, \usetikzlibrary{arrows.meta,positioning,calc,fit,backgrounds}
\begin{tikzpicture}[
    font=\small,
    box/.style={draw=DarkOrchid, fill=Thistle!55, rounded corners=3pt, align=center, minimum height=1.0cm},
    store/.style={draw=DarkOrchid, fill=white, font=\scriptsize, text width=3.95cm, minimum height=0.42cm,
                  align=center},
    flow/.style={-{Stealth[length=2.2mm]}, thick, DarkOrchid},
    lbl/.style={font=\scriptsize, align=center, fill=white, inner sep=1.5pt}
  ]
  % clients
  \node[box, minimum width=2.9cm, font=\small] (A) at (0,1.6)  {Alice (учасник)\\[-1pt]\scriptsize володіє $\mathit{GroupMasterKey}$};
  \node[box, minimum width=2.9cm, font=\small] (B) at (0,-1.6) {Bob (учасник)\\[-1pt]\scriptsize володіє $\mathit{GroupMasterKey}$};
  % server
  \node[box, minimum width=4.6cm, minimum height=4.6cm, fill=Lavender, anchor=center] (S) at (8.0,0) {};
  \node[font=\small\bfseries, anchor=north] at (S.north) {Сервер Signal};
  \node[font=\scriptsize, anchor=north, align=center] at ($(S.north)+(0,-0.45)$)
       {емітент та пруфер\\ зберігає $\mathit{ServerSecretParams}$};
  \node[store] (r1) at ($(S.center)+(0,0.05)$)
       {$\mathsf{Enc}(\mathrm{UID}_A)\;\big|\;\mathsf{Enc}(\mathrm{PK}_A)\;\big|\;\text{role}$};
  \node[store, below=0pt of r1] (r2)  {$\mathsf{Enc}(\mathrm{UID}_B)\;\big|\;\mathsf{Enc}(\mathrm{PK}_B)\;\big|\;\text{role}$};
  \node[store, below=0pt of r2] (r3)  {$\mathsf{Enc}(\text{title}),\ \mathsf{Enc}(\text{avatar}),\ \dots$};
  \node[font=\scriptsize, above=1pt of r1] {Зашифрований стан групи};
  % flows
  \draw[flow] ($(A.east)+(0,0.25)$) -- node[lbl, above=2pt, font=\scriptsize] {\textcircled{\tiny 1} daily $\mathit{AuthCredential}$\\[-1pt]
        (issued to $\mathrm{UID}_A$)} ($(S.west)+(0,1.85)$);
  \draw[flow] ($(S.west)+(0,1.25)$) -- ($(A.east)+(0,-0.25)$);
  \draw[flow] ($(A.south east)+(-0.2,0)$) -- node[lbl, pos=0.55]
        {\textcircled{\tiny 2} $\bigl(\mathsf{Enc}(\mathrm{UID}_A),\ \pi\bigr)$\\[-1pt]
        anonymous request} ($(S.west)+(0,0.2)$);
  \draw[flow] ($(S.west)+(0,-0.6)$) -- node[lbl, pos=0.5] {\textcircled{\tiny 3} encrypted entries} ($(B.north east)+(-0.2,0)$);
  \draw[flow, dashed, Orchid] ($(A.south)+(-0.4,0)$) -- node[lbl, left=1pt] {key shared\\ end-to-end} ($(B.north)+(-0.4,0)$);
\end{tikzpicture}

    \caption{Архітектура системи приватних груп в Signal (опісля ~\cite{CPZ20})}
    \label{fig:architecture}
  \end{minipage}
\end{figure}

У таблиці~\ref{tab:roles} підсумовано, хто що доводить. Важливим зауваженням є те, що сервер виконує \emph{обидві ролі} 
емітента та верифікатора: він видає \emph{AuthCredential} автентифікованому обліковому запису один раз на день, 
а згодом перевіряє анонімні представлення цього посвідчення. Тому публічна верифікованість є непотрібною, і KVAC 
досить достатньо~\cite{SignalBlog2019}.

\begin{table}[!ht]
  \centering
  \small
  \begin{tabularx}{\linewidth}{@{}l X@{}}
    \toprule
    Prover                   & Клієнт Signal учасника групи (Android, iOS, Desktop)                                    \\
    Verifier                 & Сервер Signal; Цей же сервер є емітентом повноважень                                    \\
    Твердження (membership)  & \enquote{Я маю дійсний AuthCredential на день $\tau$ для \emph{деякого} $\UID$ і $E$ це
    шифрування \emph{цього} $\UID$ за допомогою ключа \emph{цієї} групи}                                               \\
    Твердження (adding data) & \enquote{$E'$  це шифрування ключа профілю, який сервер засвідчив для цього
    $\UID$} (ProfileKeyCredential)                                                                                     \\
    Приховано & $\UID$, profile key, власне повноваження, посилання між представленнями та публікацією                 \\
    Розкрито                 & День $\tau$, відкриті параметри групи, шифротексти $E, E'$                              \\
    \bottomrule
  \end{tabularx}
  \caption{Учасники та твердження в системі приватних груп (PGS)}
  \label{tab:roles}
\end{table}

\subsection{Чому існуючих інструментів було недостатньо}

Signature-based посвідчення (CL, BBS) вирішили б проблему автентифікації, але вони є дорожчими, а їхня публічне 
підтвердження не дає жодної користі, коли емітент і верифайер збігаються~\cite{SignalBlog2019}. Існуючий KVAC від 
CMZ був ефективним, але не міг ефективно довести, що атрибут посвідчення дорівнює відкритому тексту шифрованого 
тексту~\cite{SignalBlog2019}: атрибути CMZ є скалярними величинами в $\Zq$, а шифрування 16-байтового $\UID$ як 
скалярних величин вимагало експоненціального алгоритму Ель-Гамаля над 16-бітними блоками. Прототип, побудований 
таким чином, генерував шифровані тексти розміром понад 600 байтів на один $\UID$~\cite{CPZ20}. Тому було розроблено 
новий алгебраїчний MAC, атрибути якого є елементами групи ; завдяки цьому розмір шифрованих текстів зменшився до 
64 байтів (дві стиснуті точки з групи ristretto255 простого порядку над ЕК Curve25519)~\cite{CPZ20}.

\subsection{Обсяг бібліотеки зараз}
З часу публікації оригінальної статті бібліотека доволі розширилася. Поточний модуль~\cite{libsignal} \texttt{zkgroup} 
містить, окрім облікових даних для автентифікації та профілю-ключа, резервних копій, посилань на дзвінки та 
завантаження аватарів, а також  \emph{підтвердження групових розсилок}; загальний рівень \texttt{zkcredential} 
(2023) формує такі облікові дані без необхідності щоразу писати новий доказ, а \texttt{poksho} реалізує самі докази.

\section{Опис протоколу}
\label{sec:protocol}

\subsection{Базова система доказів: \texorpdfstring{$\Sigma$}{Sigma}-протоколи за Фіатом-Шаміром}
Як ми вводили в лекції \No~3~\cite{SigmaProtocols, FS86}, що $\Sigma$-протокол для відношення $R$ -- це трикроковий 
протокол із відкритою монетою (public-coin protocol), де $a$ -- зобов'язання, $c$ -- челендж, $z$ -- відповідь, 
є \emph{повним}, \emph{special sound} (я хз як перекласти коректно), загалом: якщо існують дві прийняті стенограми 
$(a,c,z)$, $(a,c',z')$ з тим самим $a$, але різними викликами $c \neq c'$, то з них можна обчислити секрет $x$, 
а також \emph{чесний веріфайер з zero-knowledge, HVZK} -- нульове розголошення відносно чесного верифікатора, 
стенограма не несе жодної інформації про $x$. Підробка ж працює, лише якщо $c$ випадкове й не залежить від $a$, 
тобто верифікатор чесно кидає монетку. Нечесний верифікатор міг би обирати $c$ залежно від $a$, і тоді таке 
міркування не проходить. У Signal це питання знімає Фіат–Шамір: виклик рахується як геш $c = H(\dots \| a \| \dots)$, 
тому ніхто не може вибрати його "зручним"{}. Усі класичні приклади -- Schnorr~\cite{Schnorr91}, Okamoto, 
Chaum-Pedersen~\cite{CP92} є випадками одного загального протоколу: для гомоморфізму $\psi:H\to T$ доказувач 
демонструє знання про $w$, де $\psi(w)=t$ (рис.~\ref{fig:sigma}). AND-композиція отримується шляхом виконання 
підпротоколів із одним спільним запитом, який знову є доказом прообразів для гомоморфізму добутку. Перетворення 
Фіата–Шаміра замінює запит верифікатора на геш транскрипту.

Для загального протоколу кожна з трьох властивостей перевіряється в одному рядку. \emph{Повнота (completeness)}: 
$\psi(z) = \psi(r+cw) = \psi(r)\,\psi(w)^{c} = a\cdot t^{c}$. \emph{Спеціальна коректність (special soundness)}: 
два транскрипти, що приймають, з однаковим значенням $a$ дають $\psi(z-z') = t^{\,c-c'}$, отже, $w = (z-z')\cdot(c-c')^{-1}$ 
є свідченням. У групі порядку $q$ (просте число), кожна ненульова різниця $c-c'$ має обернене за модулем $q$, тому 
челенджі можна обирати з усього $\Zq$, а похибка коректності становить $1/q\approx 2^{-252}$ для алгоритму 
ristretto255. \emph{HVZK:} Каже нам обрати $z\sample H$ і $c\sample\mathcal{C}$ (н.о.р. в.в.), і обчислити $a=\psi(z)\cdot t^{-c}$; 
симульований транскрипт $(a,c,z)$ буде розподілений точно так само, як і справжній. Остання формула також 
дозволяє \texttt{poksho} передавати лише $(c,z)$ і давати перевіряючому можливість повторно порахувати $a$.

\begin{figure}[!ht]
  \centering
  \begin{minipage}{0.8\linewidth}
    \centering
    % Sigma-protocol for the preimage of a group homomorphism psi : H -> T (redrawn after Lecture 3, slides "Sigma protocol.
% Scheme" and "Generalisation of the Sigma-protocol"), with the Fiat-Shamir variant used by poksho.
% Requires: \usepackage[svgnames]{xcolor}, \usetikzlibrary{arrows.meta,positioning,calc}
\begin{tikzpicture}[
    font=\small,
    party/.style={draw=DarkOrchid, fill=Thistle!60, rounded corners=3pt, minimum width=2.6cm, minimum height=0.75cm,
                  font=\small\bfseries},
    msg/.style={-{Stealth[length=2.2mm]}, thick, DarkOrchid},
    note/.style={font=\footnotesize, align=center}
  ]
  \node[party] (P) at (0,0)   {Prover $\mathcal{P}$};
  \node[party] (V) at (7.4,0) {Verifier $\mathcal{V}$};
  \node[note, below=1mm of P] {знає $w\in H$ і $\psi(w)=t$};
  \node[note, below=1mm of V] {знає $(\psi, t)$};
  % local computations
  \node[note] at (0,-1.45)   {$r \xleftarrow{\$} H$,\; $a \gets \psi(r)$};
  \node[note] at (7.4,-2.35) {$c \xleftarrow{\$} \mathcal{C}$};
  \node[note] at (0,-3.25)   {$z \gets r + c\cdot w$};
  \node[note] at (7.4,-4.15) {приймається виключно, \\ якщо $\psi(z) = a \cdot t^{\,c}$};
  % messages
  \draw[msg] (1.7,-1.45) -- node[above, font=\footnotesize] {зобов'язання $a$} (5.7,-1.45);
  \draw[msg] (5.7,-2.35) -- node[above, font=\footnotesize] {челендж $c$} (1.7,-2.35);
  \draw[msg] (1.7,-3.25) -- node[above, font=\footnotesize] {відповідь $z$} (5.7,-3.25);
  % Fiat-Shamir box
  \node[draw=Orchid, fill=Lavender, rounded corners=3pt, align=left, font=\footnotesize, anchor=north] at (3.7,-4.75)
    {\textbf{Fiat--Shamir (як в \texttt{poksho}):}\
     $c \gets H\bigl(\text{label}\,\|\,\text{statement}\,\|\,\text{public points}\,\|\,a\,\|\,\text{msg}\bigr)$\\[1pt]
     the proof is $\pi=(c,z)$; $\mathcal{V}$ переобчислює $a=\psi(z)\cdot t^{-c}$ і далі береться геш};
\end{tikzpicture}

    \caption{$\Sigma$-протокол для прообразу гомоморфізму $\psi$~\cite{SigmaProtocols} \\
      та реалізація форми за Fiat-Shamir в \texttt{poksho}~\cite{libsignal}.}
    \label{fig:sigma}
  \end{minipage}
\end{figure}

Як можна бачити, кожен доказ у \texttt{zkgroup} є кон'юнкцією рівнянь виду $P_0 = \prod_j P_j^{\,s_j}$ між 
відкритими точками $P_j\in\G$ та таємні скалярні величини $s_j\in\Zq$; відображення $\psi(s_1,\dots,s_k) = \bigl(\prod_j P_{1j}^{s_j},\dots\bigr)$ 
є гомоморфізмом виду $\Zq^k \to \G^m$, отже, загальний протокол застосовується безпосередньо.

\begin{table}[!ht]
  \centering
  \small
  \begin{tabularx}{\linewidth}{@{}l X@{}}
    \toprule
    Арифметизація     & none; твердження є лінійними відношеннями в експоненті, відсутні схеми або системи обмежень \\
    Схема зобов'язань & Pedersen-style зобов'язання $C = G^{z}\cdot M$ щодо атрибутів елементів групи               \\
    Fiat--Shamir      & так; твердження та всі публічні точки гешуються, HMAC-SHA256-based SHO    \\
    Trusted setup     & none: усі генератори отримуються шляхом гешування фіксованих міток                          \\
    Group             & ristretto255, порядок простих чисел $q\approx 2^{252}$, 32-байтове кодування~\cite{RFC9496} \\
    Proof size        & $32\,(1+k)$ байтів для $k$ секретних скалярів: апит плюс одна відповідь на кожен скаляр     \\
    Verifier          & визначений: для верифікації потрібен секретний ключ MAC сервера (by design)                 \\
    \bottomrule
  \end{tabularx}
  \caption{Властивості системи доведення, що використовується в \texttt{zkgroup}.}
  \label{tab:properties}
\end{table}

\subsection{Алгебраїчні MACs}
Алгебраїчний MAC (Message Authentication Code) -- це MAC, тег якого є елементом групи, обчисленим за допомогою 
алгебраїчної формули низького ступеня з ключем та повідомленням, завдяки чому твердження щодо тегів можна довести 
за допомогою $\Sigma$-протоколів. Найпростішим прикладом є $\mathsf{MAC}_{\mathsf{GGM}}$~\cite{CMZ14}: з ключем 
$(x_0,x_1) \in \Zq^2$, тег для $m\in\Zq$ дорівнює $\sigma=(U,\,U^{x_0+x_1 m})$ для випадкового $U \in \G$, а під 
час верифікації перераховується друга складова. Оскільки $(U^a, U'^a)$ знову є дійсним тегом для будь-якого $a$, 
то теги можна повторно рандомізувати, що і буде основою неможливостю встановлення зв'язку.

Алгоритм MAC~\cite{CPZ20}, здійснює автентифікацію вектора елементів групи $M_1, \dots, M_n$ (за алгоритмом~\ref{alg:cpzmac}). 
Емітент оприлюднює параметри $C_W$ та $I$, які пов'язані з його ключем; згодом вони використовуються клієнтом для 
перевірки видачі та у серверній перевірці.

\begin{algorithm}[!ht]
  \small
  \caption{Алгебраїчний MAC типу CPZ над елементами групи~\cite[\S 3.1]{CPZ20}}
  \label{alg:cpzmac}
  \begin{algorithmic}[1]
    \Statex \textbf{Setup:} prime-order group $\G$ with independent generators $G_w,G_{w'},G_{x_0},G_{x_1},G_{y_1},
      \dots,G_{y_n},G_V$
    \Procedure{KeyGen}{}
    \State $w,w',x_0,x_1,y_1,\dots,y_n \sample \Zq$; \quad $W \gets G_w^{\,w}$
    \State $C_W \gets G_w^{\,w}G_{w'}^{\,w'}$; \quad
    $I \gets G_V \big/ \bigl(G_{x_0}^{x_0}G_{x_1}^{x_1}\prod_{i=1}^{n} G_{y_i}^{y_i}\bigr)$
    \State \Return $\mathit{sk}=(W,w,w',x_0,x_1,\vec y)$, \ $\mathit{iparams}=(C_W,I)$
    \EndProcedure
    \Procedure{MAC}{$\mathit{sk}, M_1,\dots,M_n$}
    \State $t \sample \Zq$, \ $U \sample \G$
    \State $V \gets W\cdot U^{x_0+x_1 t}\cdot\prod_{i=1}^{n} M_i^{\,y_i}$
    \State \Return $(t,U,V)$
    \EndProcedure
    \Procedure{Verify}{$\mathit{sk}, \vec M, (t,U,V)$}
    \State \Return $V \overset{?}{=} W\cdot U^{x_0+x_1 t}\cdot\prod_{i} M_i^{\,y_i}$
    \EndProcedure
  \end{algorithmic}
\end{algorithm}

\subsection{Веріфіковане шифрування ідентифікаторів}

Кожен $\UID$ відображається на дві точки: $M_1=\mathsf{HashToG}(\UID)$ та $M_2=\mathsf{EncodeToG}(\UID)$, оборотне 
кодування на основі алгоритму Elligator (\enquote{lizard} в коді~\cite{libsignal}). За допомогою групового ключа 
$(a_1,a_2) \in \Zq^2$ та його відкритого комітменту $A = G_{a_1}^{a_1}G_{a_2}^{a_2}$, шифротекст має вигляд:

\begin{equation}
  E_1 = M_1^{\,a_1}, \qquad E_2 = E_1^{\,a_2}\cdot M_2 .
  \label{eq:enc}
\end{equation}
При розшифруванні обчислюється $M_2 = E_2/E_1^{a_2}$, декодуєтсья $\UID$ і перевіряється \\ 
$E_1 = \mathsf{HashToG}(\UID)^{a_1}$. Така схема симетрична, детермінована та захищеною від CCA~\cite{CPZ20}: 
один і той самий UID завжди дає однаковий шифрований текст у заданій групі, що дозволяє серверу знайти запис 
користувача, який надсилає запит, за допомогою простої рівності.

\subsection{Випуск та представлення}
\label{sec:show}
На малюнку~\ref{fig:lifecycle}, який я перемалював зі статті сам власне життєвий цикл видачі. Під час видачі 
сервер обчислює MAC на атрибутах користувача ($M_1, M_2$ та redemption day) і доказує у форматі zero knowledge, 
що він використав ключ, зафіксований у $(C_W,I)$; облікові дані з профілем-ключем видаються \emph{всліпу}, на 
атрибутах, зашифрованих за алгоритмом Ель-Гамаля, тому сервер ніколи не бачить ключ профілю~\cite{CPZ20}.

\newpage % FORCED
\begin{figure}[!ht]
  \centering
  \begin{minipage}{0.95\linewidth}
    \centering
    % Life cycle of a keyed-verification anonymous credential (KVAC): issuance and unlinkable presentation.
% Requires: \usepackage[svgnames]{xcolor}, \usetikzlibrary{arrows.meta,positioning,calc}
\begin{tikzpicture}[
    font=\small,
    party/.style={draw=DarkOrchid, fill=Thistle!60, rounded corners=3pt, minimum width=3.0cm, minimum height=0.75cm,
                  font=\small\bfseries, align=center},
    msg/.style={-{Stealth[length=2.2mm]}, thick, DarkOrchid},
    lab/.style={above, font=\scriptsize, fill=white, inner sep=1pt},
    phase/.style={draw=Orchid, fill=Lavender, rounded corners=2pt, font=\scriptsize\bfseries, inner sep=2pt},
    loc/.style={font=\scriptsize, align=left}
  ]
  \node[party] (U) at (0,0)    {User, клієнт};
  \node[party] (S) at (9.0,0)  {Server\\[-2pt]\scriptsize емітент, веріфаєр};
  \draw[densely dotted, Plum] (U.south) -- (0,-8.7);
  \draw[densely dotted, Plum] (S.south) -- (9.0,-8.7);
  % ---------------- issuance
  \node[phase] at (4.5,-0.95) {Задача: автентифікація раз в день};
  \draw[msg] (0.1,-1.75) -- node[lab] {ідентифікатор $\mathrm{UID}$ (звичайний вхід)} (8.9,-1.75);
  \node[loc, anchor=east] at (8.85,-2.55) {$t\xleftarrow{\$}\mathbb{Z}_q,\ U\xleftarrow{\$}\mathbb{G}$,\\
         $V = W\,U^{x_0+x_1t}\prod_i M_i^{y_i}$};
  \draw[msg] (8.9,-3.45) -- node[lab] {$(t,U,V)$ та доказ $\pi_I$ правильності} (0.1,-3.45);
  \node[loc, anchor=west] at (0.15,-3.95) {перевірка $\pi_I$ за допомогою $(C_W, I)$};
  % ---------------- presentation
  \node[phase] at (4.5,-4.7) {Наявність: анонімно, будь-яка кількість разів};
  \node[loc, anchor=west] at (0.15,-5.5) {$z\xleftarrow{\$}\mathbb{Z}_q$:\ $C_{x_0}=G_{x_0}^zU,\ C_{x_1}=G_{x_1}^zU^t$,\\
         $C_{y_i}=G_{y_i}^zM_i,\ C_V=G_V^zV$;\ шифротекст $E$};
  \draw[msg] (0.1,-6.45) -- node[lab] {$\bigl(C_{x_0},C_{x_1},\{C_{y_i}\},C_V,E,\pi\bigr)$} (8.9,-6.45);
  \node[loc, anchor=east] at (8.85,-7.25) {$Z \gets C_V\big/\bigl(W\,C_{x_0}^{x_0}C_{x_1}^{x_1}\prod_i C_{y_i}^{y_i}\bigr)$,\\
         прийняти виключно, якщо $\pi$ підтверджено обраним $Z$};
  \node[loc, anchor=west, DarkOrchid, text width=7.6cm] at (0.15,-8.2)
       {Оскільки щоразу нове $z$ використовується $\Rightarrow$ презентації не можна пов'язати між собою та з моментом видачі};
\end{tikzpicture}

    \caption{Випуск та представлення посвідчення з верифікацією за ключем у \texttt{zkgroup} (after~\cite{CPZ20}).}
    \label{fig:lifecycle}
  \end{minipage}
\end{figure}

\noindent Для представлення клієнт обирає нове $z\sample\Zq$ Z і замінює кожен компонент Pedersen-style зобов'язанням: 
$C_{x_0}=G_{x_0}^{z}U$, $C_{x_1}=G_{x_1}^{z}U^{t}$, $C_{y_i}=G_{y_i}^{z}M_i$, $C_V=G_V^{z}V$. Використосуючи свій 
секретний ключ, сервер обчислює:
\begin{equation}
  Z = \frac{C_V}{W\,C_{x_0}^{x_0}\,C_{x_1}^{x_1}\prod_i C_{y_i}^{y_i}}
  = \frac{G_V^{z}}{\bigl(G_{x_0}^{x_0}G_{x_1}^{x_1}\prod_i G_{y_i}^{y_i}\bigr)^{z}}
  \cdot\underbrace{\frac{V}{W\,U^{x_0+x_1t}\prod_i M_i^{y_i}}}_{=\,1 \text{ for a valid tag}}
  = I^{z},
  \label{eq:Z}
\end{equation}
отже, дійсні облікові дані еквівалентні твердженню $Z=I^z$ для $z$, відомого клієнту. Для AuthCredential клієнт 
доводить за допомогою одного Fiat-Shamir $\Sigma$-доказу~\cite[\S 5]{CPZ20},
\begin{equation}
  \pi_A = \mathrm{PK}\left\{ (z, z_0, z_1, t, a_1, a_2) :
  \begin{array}{@{}l@{\qquad}l@{}}
    Z = I^{z}                                           & \text{валідний MAC}              \\
    C_{x_1} = C_{x_0}^{\,t}\,G_{x_0}^{z_0}\,G_{x_1}^{z} & \text{послідовні } t \ (z_0=-tz) \\
    A = G_{a_1}^{a_1} G_{a_2}^{a_2}                     & \text{зареестрований ключ групи} \\
    E_1 = C_{y_1}^{\,a_1}\,G_{y_1}^{z_1}                & E_1 = M_1^{a_1} \ (z_1=-za_1)    \\
    C_{y_2}/E_2 = G_{y_2}^{z}\big/E_1^{\,a_2}           & E_2 = E_1^{a_2}M_2
  \end{array}
  \right\}.
  \label{eq:piA}
\end{equation}

Перший рядок -- це перевірка повноважень за рівнянням~\ref{eq:Z}; решта рядків це \emph{verifiable encryption} 
шифрування, які пов'язують сертифіковані атрибути $M_1, M_2$ із шифрованим текстом за рівнянням~\ref{eq:enc}. 
Кожен рядок є лінійним щодо таємних експонентів, тому рівняння~\ref{eq:piA} є окремим доказом прообразу гомоморфізма, 
тобто AND-композицією (як ми вводили це в лекції \No~3~\cite{SigmaProtocols}). Розкритий "день погашення"{} 
обробляється додатковим рядком такої самої форми~\cite{CPZ20}.

Я ще вирішив в звіт додати ще іграшковий приклад, на якому я розбирався з штучним інтелектом, як то всьо працює:
\begin{examplebox}[Приклад на алгебрі в групі порядку 11]
  Розглянемо підгрупу порядку $q=11$ над полем $\mathbb{Z}_{23}^{\times}$ породжену генератором $g=2$.

  \emph{MAC та re-randomisation.} Оберемо ключ $(x_0,x_1)=(3,5)$, повідомлення $m=2$, $U=g^2=4$: $x_0+x_1m = 13 \equiv 2 \pmod{11}$, 
  тож тег дорівнюватиме $(U,U')=(4,\,4^2=16)$. Піднесення обох частин до $a=4$ дає $(4^4,16^4)=(3,9)$, і справді 
  $3^2=9$: тег, що виглядає новим, для того самого $m$, який емітент не може пов'язати з $(4, 16)$ (unlinkability).

  P.S. Перерандомізація ховає, \emph{який саме} тег показують, але \emph{не ховає} повідомлення $m$. Видавець знає 
  ключ і може перевірити, чи $9 = 3^{x_0 + x_1 m}$, перебираючи $m$. Тому в справжньому протоколі, як зазначалося 
  вище клієнт не надсилає $U, V$ і атрибути відкрито. Замість цього він надсилає зобов'язання зі новим $z$. 
  Атрибути сховані, а коректність доводиться $\Sigma$-протоколом ($Z = I^z$).

  \emph{Доказ Шнорра та Фіата–Шаміра}. Таємний ключ $x = 7$, відкритий ключ $X = g^7 = 13$. Прувер обирає $r = 3$, 
  і надсилає $a = g^3 = 8$; нехай результат гешування $c = 4$; тоді $z = r + cx = 31 \equiv 9$. Верифікатор перевіряє 
  $g^9=6$ та $a \cdot X^c = 8 \cdot 13^4 = 8 \cdot 18 \equiv 6 \pmod{23}$. Якщо доказувач повторно використав би 
  $r=3$ для другого челенджу $c'=9$ ($z'=0$), то будь-хто міг би витягти $x=(z-z')/(c-c')=9\cdot 6^{-1}=9\cdot 2\equiv 7$.
  Це власне той екстрактор зі спеціальною коректністю, про який йшлося на лекції, перетворюється на атаку. Через 
  це \texttt{poksho} виводить nonces на основі свіжої випадковісті \emph{та} гешу свідчення й повідомлення.
\end{examplebox}

\subsection{Модифікації та особливості реалізації}
Розгорнутий протокол відрізняється кількома особливостями:
\begin{enumerate}
  \item \textbf{New MAC.} Атрибути є елементами групи, а не скалярними величинами (алгоритм~\ref{alg:cpzmac}); 
    саме це робить можливими 64-байтові підтвержувані шифровані тексти~\cite{CPZ20}.
  \item \textbf{Strong Fiat--Shamir.} \texttt{poksho} приймає мітку протоколу, канонічний опис твердження, усі 
    відкриті точки, зобов'язання та optional повідомлення перед стисненням запиту~\cite{libsignal}. Відсутність 
    statement призвело б до \enquote{слабкого} (weak) варіанту, який, як відомо, є некоректним (незахищеним) 
    для адаптивно обраних зломвисником тверджень~\cite{BPW12}.
  \item \textbf{Synthetic nonces.} Випадковість зобов'язань походить від свіжої випадковості та хешу свідчення 
    \& повідомлення, тому "бракований"{} генератор випадкових чисел не може спричинити nonce reuse, як у ілюстративному 
    прикладі~\cite{libsignal}.
  \item \textbf{Self-check and compact proofs.} Доказувач перевіряє власний доказ перед його оприлюдненням (захист
    від fault attacks); докази кодуються у вигляді $(c, z_1, \dots, z_k)$; верифікатор перераховує зобов'язання 
    як $a=\psi(z)\cdot t^{-c}$~\cite{libsignal}
  \item \textbf{Prime-order group usage.} Група простого порядку ristretto255~\cite{RFC9496} допомагає усунути 
    значну частину криптографічних вразливостей Curve25519 (класичне джерело прихованих помилок у доказах над ЕК).
\end{enumerate}

\section{Аналіз}
\label{sec:analysis}

\subsection{Криптографічні припущення}
Система ґрунтується на чотирьох припущеннях~\cite{CPZ20}. Використання \emph{DDH} в $\G$, для забезпечення безпеки 
нового MAC та шифрування. (ii) Непідробність $\mathsf{MAC}_{\mathsf{GGM}}$, для якої єдині відомі докази належать 
до \emph{моделі загальної групи}; автори зводять новий MAC до неї та до DDH (сильна непідробність у разі chosen-message 
та verification attacks в моделі випадкового оракула). (iii) \emph{Модель випадкового оракула} для Fiat-Shamir 
та $\mathsf{HashToG}$. (iv) Припущення DL, а також властивість \emph{strong extraction} системи доказів.

Для порівняння, безпека BBS доведена за припущенням $q$-SDH у групах спарювання~\cite{TZ23}, параметризованому 
припущенні, яке є сильнішим за DDH, тоді як залежність від GGM є головною теоретичною слабкою стороною конструкції 
Signal.

\subsection{Відомі вразливості і фейли}
Скільки шо я, що claude, gemini з deepseek не копали, звітів про вразливості, що стосуються саме \texttt{zkgroup} 
або \texttt{poksho} не знайшли. Логічно, що це не є гарантією безпеки. Але відомі типи помилок у технологіях, з 
яких воно зібрано, і відповідні класи ризиків добре задокументовані в інших джерелах:

\begin{itemize}
  \item \textbf{Weak Fiat--Shamir.} У 2022 році компанія Trail of Bits розкрила вразливості \enquote{Frozen Heart} 
    у реалізаціях Girault's proof, Bulletproofs and PlonK, спричинені недостатнім гешуванням транскрипту~\cite{FrozenHeart22}; 
    раніше цей слабкий варіант порушив аналіз безпеки системи голосування Helios~\cite{BPW12}. Як результат 
    \texttt{poksho} використовує сильний Фіата–Шаміра (описував в розділі~\ref{sec:protocol}).
  \item \textbf{Повторне використання нонсів / погана рандомізація} призводить до витоку свідчення(witness); це 
    пом'якшується за допомогою використання синтетичних нонсів.
  \item \textbf{Модель довіри для верифікації з використанням ключа.} Той, хто володіє серверним ключем, може 
    створювати облікові дані для будь-якого $\UID$. Це є прийнятним лише тому, що той самий сервер і забезпечує 
    контроль доступу; це не дозволяє серверу розшифрувати стан групи, для чого потрібен вже груповий ключ.
  \item \textbf{Детерміноване шифрування} розкриває серверу рівність записів \emph{у межах групи} -- це передбачено, 
    але про це слід пам'ятати, оскільки можливо отримати витік інформації.
  \item \textbf{Складність.} Шифрування, MAC та докази поєднуються у спосіб, що не відповідає принципу "чорного 
    ящика"{}(non-black-box), що ускладнює аудит системи та її перенесення~\cite{Connell26}.
\end{itemize}

\subsection{Перспективи на постквантову епоху}
Кожен компонент у алгоритмі ristretto255 ґрунтується на дискретних логарифмах і може бути зламаний за допомогою 
алгоритму Шора. Найбільш нагальним наслідком є широкозастосований сценарій \emph{harvest-now-decrypt-later, HNDL}, 
що полягав в збиранні шифротекстів, що містять конфіденційну інформацію про приналежність до групи, зберігаються 
сьогодні, і можуть бути розшифровані майбутнім квантовим зловмисником. Коннелл зі співавторами~\cite{Connell26} 
стверджують, що проста заміна на квантово-безпечні примітиви навряд чи дасть ефективний протокол, оскільки 
квантово-безпечні повноваження та верифіковуване шифрування потребують складних доказів; вони пропонують модульне 
перепроектування на основі підписів із можливістю повторної рандомізації ключів, що дозволяє перенести першими 
HNDL-вразливі частини

\subsection{Продуктивність}
\label{sec:perf}
В табличці~\ref{tab:bench} я зібрав опубліковані результати тестування і власні запуски алгоритму на Rust. 
Комплілював \texttt{libsignal} на коміті \texttt{8509425}~\cite{libsignal}, запустив стандартні вбудовані тести 
(benchmarks) та невелику програмку (яку написав Claude), що виводить розміри серіалізованих даних.

Усі вимірювання я проводив на своєму ноутбуці з процесором Intel Core i7-1065G7 і Windows. У репозиторії зазначалося, 
що версія набору інструментів Rust зафіксована у файлі \texttt{rust-toolchain}, тому використовувалася саме вона. 
Єдиною додатково встановленою залежністю був компілятор Protocol Buffers \texttt{protoc}. Запустив команду: \\
\texttt{cargo bench -p zkgroup --bench zkgroup\_benchmarks}, яка використовує фреймворк \texttt{criterion}: 
кожна операція спочатку виконується протягом 3 секунд для "розігріву"{}, а потім повторюється 100 разів, і 
\texttt{criterion} видає точкову оцінку часу на одну операцію разом із 95\,\% довірчим інтервалом. У таблиці 
нижче наведено точкові оцінки, довірчі інтервали для всіх операцій, які не перевищували $\pm 3\,\%$.

Розміри серіалізованих даних не залежать від апаратного забезпечення, тому вони вимірювалися за допомогою окремої 
короткої програми (\texttt{sizes.rs}, розміщеної в \texttt{rust/zkgroup/examples}), яка створює кожен об'єкт із 
фіксованими тестовими початковими значеннями, серіалізує його за допомогою власної функції бібліотеки 
\texttt{zkgroup::serialize} та виводить довжину в байтах. Для порівняння з BBS було штучно зроблено невееликий 
"проєкт"{} Rust, який проводить тестування продуктивності реалізації \texttt{zkryptium} (версія 0.8.1, набір 
шифрів BLS12-381-SHA-256) з тими самими налаштуваннями \texttt{criterion} (3 секунди розігріву + 100 повторів) 
на тій самій машині. Його параметри це чотири підписані повідомлення, одне відкрите та три приховані, були обрані 
саме так, щоб відображати атрибути AuthCredential.

\begin{table}[!ht]
  \centering
  \small
  \caption{Розміри та час виконання основних об'єктів \texttt{zkgroup}. Опубліковані значення: Intel Core i7-8650U~\cite{CPZ20};
    власні значення: Intel Core i7-1065G7 (laptop, Windows), точкові оцінки, наведені в \texttt{criterion}}
  \label{tab:bench}
  \setlength{\tabcolsep}{4.5pt}
  \begin{tabular}{@{}l rrr rrr@{}}
    \toprule
                                            & \multicolumn{3}{c}{Опубліковані~\cite{CPZ20}} & \multicolumn{3}{c}{Власні виміри}                                               \\
    \cmidrule(lr){2-4}\cmidrule(l){5-7}
    Object                                  & bytes                                      & create, ms                          & verify, ms & bytes & create, ms & verify, ms \\
    \midrule
    UidCiphertext                           & 64                                         & 0.13                                & 0.18       & 65    & 0.085      & 0.158      \\
    ProfileKeyCiphertext                    & 64                                         & 0.13                                & 0.86       & 65    & 0.077      & 0.550      \\
    AuthCredential response$^{a}$           & 361                                        & 1.95                                & 0.98       & 425   & 1.15       & 0.51       \\
    AuthCredential presentation$^{a}$       & 493                                        & 2.16                                & 1.17       & 633   & 1.85       & 0.95       \\
    ProfileKeyCredential request            & 329                                        & 1.48                                & --         & 329   & 1.01       & --         \\
    ProfileKeyCredential response$^{b}$     & 457                                        & 2.70                                & 0.97       & 497   & 2.01       & 0.69       \\
    ProfileKeyCredential presentation$^{b}$ & 713                                        & 2.87                                & 1.53       & 721   & 1.97       & 1.12       \\
    \bottomrule
    \multicolumn{7}{@{}p{0.97\linewidth}@{}}{\footnotesize $^{a}$ варіант, \texttt{AuthCredentialWithPniZkc}
    який в продакшині, сертифікує два ідентифікатори (ACI та PNI) і побудований на основі \texttt{zkcredential}, 
    і тому має більший розмір. \newline $^{b}$ сучасний \texttt{Expiring} варіант містить час закінчення терміну дії. 
    Розміри серіалізованих даних містять один байт версії.}
  \end{tabular}
\end{table}

Усі операції займають від однієї до двох мілісекунд, а кожне повідомлення вміщується значно менше ніж в один 
кілобайт. Завантаження групи з $n$ членів, згідно з~\cite{CPZ20}, вимагає приблизно $128n$ байтів і $n(0.86+0.18)$\,ms 
на дешифрування; Згідно моїх показників вийшло $\approx 0.71$\,ms на кожного члена, тобто менше ніж 0.1\,с.

\subsection{Чи можна заміняти протокол і якою ціною?}
\label{sec:alternatives}
\paragraph{BBS signatures.} Очевидним кандидатом є схема BBS, описана в CFRG\cite{BBSdraft}, аналіз безпеки якої 
базується на роботі Тессаро та Чжу~\cite{TZ23}. Вона узагальнено зображена в алгоритмі~\ref{alg:bbs}:

\newpage %Forced
\begin{algorithm}[!ht]
  \caption{BBS signatures над ЕК BLS12-381 (спрощений вигляд з~\cite{BBSdraft})}
  \label{alg:bbs}
  \small
  \begin{algorithmic}[1]
    \Statex \textbf{Setup:} pairing $h:\G_1\times\G_2\to\G_T$; generators $P_1, Q_1, H_1,\dots,H_L\in\G_1$,
    $\mathit{BP}_2\in\G_2$
    \Procedure{KeyGen}{}
    \State $\mathit{SK}\sample\Zq$; \quad $W \gets \mathit{SK}\cdot\mathit{BP}_2$ \Comment{public key, 96 bytes}
    \EndProcedure
    \Procedure{Sign}{$\mathit{SK}, m_1,\dots,m_L$}
    \State $B \gets P_1 + Q_1\cdot\mathit{domain} + \sum_{i} H_i\cdot m_i$; \quad
    $e \gets \mathsf{hash\_to\_scalar}(\mathit{SK}, \vec m, \mathit{domain})$
    \State $A \gets B\cdot (\mathit{SK}+e)^{-1}$; \quad \Return $(A,e)$ \Comment{80 bytes}
    \EndProcedure
    \Procedure{Verify}{$W, \vec m, (A,e)$}
    \State \Return $h(A, W)\cdot h(A\cdot e - B, \mathit{BP}_2) \overset{?}{=} 1_{\G_T}$
    \EndProcedure
    \Procedure{ProofGen}{$W, (A,e), \vec m, \text{disclosed set}$}
    \State randomise $A$ into $(\bar A, \bar B, D)$ and prove with a $\Sigma$-protocol knowledge of $e$, the
    randomness and the $U$ undisclosed $m_j$; output $3\cdot 48 + (4+U)\cdot 32$ bytes
    \EndProcedure
  \end{algorithmic}
\end{algorithm}

Щоб порівняти витрати в рівних умовах, було проведено тестування продуктивності реалізації з відкритим кодом 
\texttt{zkryptium}~\cite{zkryptium} (набір шифрів BLS12-381-SHA-256) із $L=4$ повідомленнями, з яких $U=3$ залишаються 
прихованими, що відтворює AuthCredential (рис.~\ref{fig:bench}). Було обрано саме її оскільки вона реалізує поточний 
проект CFRG і компілюється як бібліотека, написана суто на Rust (мені для Вінди це було в самий раз). На відміну 
від сильно оптимізованої арифметики кривих, що використовується в \texttt{libsignal} (не покладається на 
оптимізований на асемблері код спарювання) Тому, порівняння є справедливим щодо апаратного забезпечення, 
але не щодо "зрілості"{} реалізації.

\begin{figure}[!ht]
  \centering
  \begin{minipage}{0.75\linewidth}
    \centering
    % Own measurements (criterion, Intel Core i7-1065G7 @ 1.30 GHz, Windows laptop, 2026-11-08):
% libsignal zkgroup AuthCredentialWithPniZkc vs. BBS (zkryptium 0.8.1, BLS12-381-SHA-256, L = 4, U = 3).
% Requires: \usepackage{pgfplots}, \pgfplotsset{compat=1.18}, \usepackage[svgnames]{xcolor}
\begin{tikzpicture}
  \begin{axis}[
      ybar, bar width=13pt, width=11cm, height=6.0cm,
      ymin=0, ymax=15.5, ylabel={Час, ms}, ylabel style={font=\small},
      symbolic x coords={Issue/Sign, Present/ProofGen, Verify},
      xtick=data, xticklabel style={font=\small},
      yticklabel style={font=\footnotesize},
      enlarge x limits=0.22,
      legend style={at={(0.5,1.03)}, anchor=south, font=\footnotesize, draw=Plum, legend columns=1},
      nodes near coords, nodes near coords style={font=\scriptsize, /pgf/number format/.cd, fixed, precision=2},
      ymajorgrids, grid style={Thistle!70}
    ]
    \addplot[draw=DarkOrchid, fill=Orchid!70] coordinates
      {(Issue/Sign,1.15) (Present/ProofGen,1.85) (Verify,0.95)};
    \addplot[draw=Plum!80!black, fill=Thistle] coordinates
      {(Issue/Sign,7.04) (Present/ProofGen,13.05) (Verify,12.42)};
    \legend{{zkgroup AuthCredential (ristretto255, KVAC)},{BBS, $L=4$, $U=3$ (BLS12-381, pairings)}}
  \end{axis}
\end{tikzpicture}

    \caption{Власні виміри: \texttt{zkgroup} AuthCredential vs.\ BBS ($L=4$, $U=3$).}
    \label{fig:bench}
  \end{minipage}
\end{figure}
Часові параметри BBS значною мірою залежать від бібліотеки спарювання. \texttt{zkgroup} виконує лише множення 
точок у \texttt{ristretto255}. BBS для кожної перевірки обчислює спарювання (pairing) на \texttt{BLS12-381}, а 
це за своєю природою значно дорожча операція, ніж множення точки на еліптичній кривій. На малюнку стовпчики 
показують \textit{порядок} величини, а не точне співвідношення.

Заміна \texttt{zkgroup} на BBS потребувала б таких витрат: (a) спарювання -- перевірка відтворення з тестованою 
нами реалізацією виконувалася на порядок повільніше; оптимізована бібліотека спарювання дозволила б зменшити цей 
розрив, але спарювання все одно залишається набагато ресурсомісткішим, ніж арифметика ristretto255, що має важливе 
значення для сервера, який перевіряє запити від мільйонів клієнтів; (b) другої еліптичної кривої (BLS12-381) у 
кодовій базі, побудованій на Curve25519; (c) І найголовніше, оскільки повідомлення BBS є скалярами, то для 
верифікованого шифрування з рівняння~\ref{eq:enc} знадобилися або докази для двох різних груп, або перенесення 
всього шару шифрування на $\G_1$ з BLS12-381, що знову створювало б саме ту проблему, яка стала приводом створення 
CPZ MAC. Перевагами були б публічна перевірюваність, яка Signal ой як не потрібна. Варто завважити, що звичайний 
доказ BBS засвідчує менше, ніж презентація \texttt{zkgroup}, яка додатково доводить правильність шифрованого тексту.

\textbf{Шо по інших варіантах?} В таблиці~\ref{tab:alternatives} згрупував інших кандидатів. BBS-MAC~\cite{BBDT16} 
та його оптимізація $\mu$BBS~\cite{Orru24} зберігають keyed verification і уникають спарювання, залишаючись 
сумісними з проектом BBS; однак він же і зазначає, що $\mu$CMZ (оптимізований CMZ, зменшено накладні 
витрати на видачу облікових даних) є ефективнішим під час видачі, тоді як презентація є менш ефективною для більш 
ніж одного атрибута~\cite{Orru24}, і їхні атрибути знову ж таки є скалярними. SNARK загального призначення міг би 
перевіряти MAC всередині схеми, але це додало б арифметизацію, часто це надійне налаштування, та значно більші 
витрати стороні прувера (а це наш телефон) для тверджень, які й так є лінійними (в показнику степенів).

\begin{table}[!ht]
  \centering
  \small
  \begin{tabularx}{\linewidth}{@{}l X X@{}}
    \toprule
    Альтернатива                             & Переваги                                                                      & Витрати/недоліки \\
    \midrule
    BBS~\cite{BBSdraft,TZ23}                 & стандартний канал, можливість публічної перевірки, доказ $q$-SDH без GGM
                                             & спарювання, потреба в другій ЕК, скалярні атрибути ламають дешевий encryption                    \\
    BBS-MAC / $\mu$BBS~\cite{BBDT16,Orru24}  & без спарювань, BBS-сумісна, сучасні поняття безпеки
                                             & скалярні атрибути; нова модель (2024)                                                            \\
    $\mu$CMZ~\cite{Orru24}, ARC~\cite{ARC25} & заходи з стандартизації для KVAC
                                             & необхідність CPZ-style розширення для зашифрованих атрибутів                                     \\
    Generic SNARK                            & можливість довести майже будь яке обчислення (expressiveness)
                                             & значно повільніша перевірка                                                                      \\
    Quantum-safe redesign~\cite{Connell26}   & стійкість до HNDL і алгоритму Шора
                                             & нова архітектура, яка все ще досліджується                                                       \\
    \bottomrule
  \end{tabularx}
  \caption{Можливі заміни credential рівня в \texttt{zkgroup}}
  \label{tab:alternatives}
\end{table}

\section*{Висновки}
\phantomsection
\addcontentsline{toc}{section}{Висновки}

Приватні групи Signal демонструють, як zero-knowledge додаток будується навколо конкретної моделі довіри. Оскільки 
сервер є одночасно емітентом і верифікатором, розробники обрали посвідчення з верифікацією за ключем, а щоб 
зробити верифіковане шифрування ідентифікаторів недорогим, вони запровадили використання алгебраїчного MAC на 
елементах групи. Кожен доказ є $\Sigma$--протоколом для прообразу гомоморфізму AND-композицією відношень типів 
Schnorr-, Okamoto- і Chaum-Pedersen-type, перетворений на неінтерактивне доведення за допомогою сильного 
перетворення Фіата–Шаміра без арифметизації та без trusted налаштування. Зроблені мною вимірювання підтверджують 
опубліковану ефективність.

Основними слабкими сторонами є залежність від загальної моделі групи для базового MAC, складність композиції, що 
не є "чорною скринькою"{}, і, перш за все, відсутність квантової стійкості з пов'язаним із цим ризиком 
harvest-now-decrypt-later. BBS, хоча й стандартизований та має публічну перевірку, став би кроком назад для цього 
конкретного застосування: він вимагає спарювань, у нашій конфігурації його перевірка була в рази повільнішою, а 
його скалярні повідомлення суперечать шифруванню за елементами групи, що лежить в основі системи. Отже, реалістичним 
розвитком є не перехід на BBS, а/або новіші схеми верифікації з ключем ($\mu$CM, стандартизовані $\Sigma$-протоколи), 
або квантово-безпечна переробка групової системи в цілому.
